\subsection{Declaración de Alcance del Proyecto}

Este proyecto se ejecuta sobre la plataforma Paralegales, un producto en operación con base de usuarios activa, y contempla una intervención técnica orientada a mitigar la deuda arquitectónica acumulada durante su etapa inicial de desarrollo sin comprometer la continuidad del servicio en producción. La estrategia adoptada es una \textbf{evolución arquitectónica híbrida} que combina Firebase y AWS de manera deliberada, en lugar de una migración total de proveedor de nube. Esta decisión condiciona de manera explícita el espacio de acción del trabajo y delimita lo que está dentro y fuera de su alcance.

Dentro del alcance se encuentran las cuatro líneas de intervención que materializan la estrategia descrita: (i) la reubicación de la lógica de negocio crítica del cliente hacia entornos de ejecución del lado del servidor; (ii) la refactorización de la arquitectura interna del \textit{frontend} y del \textit{backend} bajo patrones modulares que aíslen responsabilidades y permitan la evolución independiente de cada componente; (iii) la construcción de un sistema transversal de observabilidad estructurada con su correspondiente infraestructura como código; y (iv) la consolidación de prácticas automatizadas de aseguramiento de calidad continuo a lo largo del ciclo de vida del desarrollo. Estas líneas se aplican sobre los módulos funcionales activos de la plataforma —autenticación, gestión de procesos judiciales y comunicación masiva mediante campañas—, que constituyen la materia sobre la cual se materializan las decisiones técnicas adoptadas.

Quedan fuera del alcance de este trabajo: (a) la sustitución de Firebase Authentication por un proveedor alternativo de gestión de identidades; (b) la migración del motor de persistencia desde Cloud Firestore hacia otro sistema de bases de datos —incluida la migración hacia Amazon DynamoDB que fue evaluada y descartada por la organización durante la fase exploratoria—; (c) la sustitución total del proveedor de nube, es decir, la transición a un entorno 100\% AWS; y (d) la incorporación de funcionalidades de negocio nuevas no relacionadas con los ejes de la intervención arquitectónica. Esta delimitación es coherente con la hipótesis del trabajo: el valor esperado se obtiene mediante la reconfiguración de las relaciones entre los subsistemas existentes y la introducción de capacidades transversales, no mediante la sustitución disruptiva del sustrato tecnológico que actualmente soporta el producto.
